\documentclass[11pt]{article}

% Plain article style for self-study reports that do not need a conference
% layout. Use the neurips style for publication-shaped delegated reports.

\usepackage[margin=1in]{geometry}
\usepackage[T1]{fontenc}
\usepackage{amsmath,amssymb}
\usepackage{booktabs}
\usepackage{graphicx}
\usepackage[colorlinks=true,linkcolor=blue,urlcolor=blue,citecolor=blue]{hyperref}
\usepackage[numbers,sort&compress]{natbib}

\title{Measuring the $1/\sqrt{d_k}$ Scale in Attention Logits}
\author{self-study-with-ai}
\date{\today}

\begin{document}
\maketitle

\begin{abstract}
The $1/\sqrt{d_k}$ factor in scaled dot-product attention is justified by a
variance argument that its source states but never measures. This report
measures it under exactly the assumptions the argument makes. Unscaled logit
variance tracks $d_k$ across two orders of magnitude, and the scale holds
variance within about one percent of unity. Softmax concentration over 64
competing keys climbs from a mean largest probability of $0.4215$ to $0.9290$
without the scale, and holds near $0.106$ with it, which is still about seven
times the uniform $0.0156$. Every figure comes from a recorded run with an
independent replication behind it. Nothing here measures gradients, so the
mechanism the source suspects remains untested.
\end{abstract}

\section{Introduction}

A companion study established what the source argues: for queries and keys
with independent, zero-mean, unit-variance components, a $d_k$-dimensional dot
product has variance $d_k$, so dividing by $\sqrt{d_k}$ restores unit
variance. That is algebra, and the source offers no measurement of it. This
report supplies one. The question is narrow on purpose: under the assumptions
the argument makes, do the numbers behave as the algebra says, and how far
from the idealized values do they actually sit?

\section{Background}

Scaled dot-product attention scores a query against each key by an inner
product, divides by $\sqrt{d_k}$, and normalizes the scores with a softmax.
The variance argument is conditional: it assumes independence and unit
variance in the components. This study imposes those assumptions in its
generator rather than testing them, so what follows is a check on the algebra,
not on a trained model.

\section{Methods and Sources}

This study cites no external sources. Its evidence is its own runs, recorded
in the study's dossier with per-run manifests, content hashes, and a claims
ledger.

For each $d_k \in \{8, 32, 128, 512\}$ the experiment samples 20000 independent
query and key pairs from a standard normal, records the population variance of
the resulting logits with and without the scale, and separately samples 2000
softmax rows of 64 competing keys to record the mean largest probability. The
scaled and unscaled figures come from the same sampled dot products, so the
comparison carries no sampling confound.

Two full runs were recorded: a main run at seed 20260901 and an independent
replication at seed 771244. Figures below come from the main run; the
replication is cited wherever the spread between runs is material. The design
was frozen in the experiment plan before either run.

\section{Findings}

\subsection{Unscaled variance grows linearly in the key dimension}

Table~\ref{tab:variance} gives the measured variances. The unscaled column
tracks $d_k$: relative to $d_k = 8$, the measured ratios are $4.03$, $15.80$,
and $63.18$ against the predicted $4$, $16$, and $64$.

\begin{table}[t]
\centering
\caption{Measured logit variance, main run at seed 20260901, 20000 pairs per dimension.}
\label{tab:variance}
\begin{tabular}{rrrr}
\toprule
$d_k$ & unscaled & scaled & scaled, replication \\
\midrule
8 & 8.0171 & 1.0021 & 1.0259 \\
32 & 32.273 & 1.0085 & 0.9927 \\
128 & 126.703 & 0.9899 & 0.9692 \\
512 & 506.554 & 0.9894 & 0.9990 \\
\bottomrule
\end{tabular}
\end{table}

\subsection{The scale holds variance near unity}

The scaled column of Table~\ref{tab:variance} stays within $1.1$ percent of
unity at every dimension in the main run. The replication column is the honest
resolution of this design: it stays within $3.1$ percent, with its largest
deviation at $d_k = 128$. Reporting the scale as producing exactly unit
variance would overstate what two runs of this size can show.

\subsection{Concentration is bounded, not removed}

Table~\ref{tab:softmax} gives the mean largest softmax probability over 64
competing keys. Without the scale it climbs steeply with $d_k$, which is the
behavior the scale exists to prevent. With the scale it is flat across
dimensions at roughly $0.106$.

\begin{table}[t]
\centering
\caption{Mean largest softmax probability over 64 keys, main run, 2000 rows per dimension. Uniform is 0.0156.}
\label{tab:softmax}
\begin{tabular}{rrr}
\toprule
$d_k$ & unscaled & scaled \\
\midrule
8 & 0.4215 & 0.1045 \\
32 & 0.7083 & 0.1071 \\
128 & 0.8497 & 0.1061 \\
512 & 0.9290 & 0.1077 \\
\bottomrule
\end{tabular}
\end{table}

That flat value deserves emphasis, because it is easy to misread. At $0.106$
the scaled softmax still puts about seven times the uniform mass on its
largest entry. The scale stops concentration from growing with dimension. It
does not produce a diffuse distribution, and a reader who takes the flat
column as evidence of near-uniform attention has drawn a conclusion these
numbers do not support.

\section{Related Work}

Out of scope by the brief. This report measures one stated argument and makes
no survey claim about later normalization proposals.

\section{Limitations}

The generator imposes the independence and unit variance the argument assumes,
so nothing here transfers to trained attention, whose queries and keys satisfy
neither. Sampling error is visible in the third significant figure, and the
two runs differ by up to $2.4$ percent on unscaled variance. Most importantly,
no gradient was measured. The source suspects that unscaled logits push the
softmax into a small-gradient region; this study measures concentration, which
is consistent with that story but does not establish it. Settling it needs a
gradient measurement, which is a separate study.

\section{Conclusion}

Under the assumptions the argument makes, the algebra holds up: unscaled logit
variance grows with $d_k$, the $1/\sqrt{d_k}$ scale returns it to roughly
unity, and softmax concentration stops growing with dimension. What the
measurement adds beyond the algebra is calibration, namely how far the real
numbers sit from the idealized ones and how much of the concentration story
the scale actually addresses.

\end{document}
